%% LyX 2.4.4 created this file.  For more info, see https://www.lyx.org/.
%% Do not edit unless you really know what you are doing.
\documentclass[english]{article}
\usepackage[T1]{fontenc}
\usepackage[latin9]{inputenc}
\usepackage{units}
\usepackage{enumitem}
\usepackage{amsmath}
\usepackage{amsthm}
\usepackage{amssymb}
\usepackage{geometry}
\geometry{verbose,tmargin=1in,bmargin=1in,lmargin=1in,rmargin=1in}

\makeatletter
%%%%%%%%%%%%%%%%%%%%%%%%%%%%%% Textclass specific LaTeX commands.
\numberwithin{equation}{section}
\numberwithin{figure}{section}
\newlength{\lyxlabelwidth}      % auxiliary length 

%%%%%%%%%%%%%%%%%%%%%%%%%%%%%% User specified LaTeX commands.
\usepackage{fancyhdr}
\usepackage{lastpage}
\pagestyle{fancy}
\usepackage{enumitem}
\usepackage{ifthen}
\everymath=\expandafter{\the\everymath\displaystyle}
%\DeclareMathSizes{10}{10}{10}{10}
\usepackage{multicol}

\usepackage{tikz}
\usepackage{tikz-3dplot}
\usetikzlibrary{calc,arrows}

\usetikzlibrary{patterns}
\usetikzlibrary{plotmarks}
\usepackage{pgfplots}
\pgfplotsset{compat=newest}

\usepackage{calc}
\usepackage[nomessages]{fp}

\usepackage{datetime}
% Added by lyx2lyx
\setlength{\parskip}{\smallskipamount}
\setlength{\parindent}{0pt}

\makeatother

\usepackage{babel}
\begin{document}
\renewcommand{\author}{Dr.\,James Ashe}

\newcommand{\classNum}{336}

\newcommand{\classSec}{A}

\newcommand{\examType}{HW 6.1}

\renewcommand{\date}{\formatdate{20}{11}{2013}}

%%First page's header%%%%%%%%%%%%%%%%%%%%%%
\fancypagestyle{firststyle} %%header settings for first pagr
{
	\fancyhead[L]{MTH \classNum{}(\classSec{}) \author{}\\
	\textbf{\examType{}} ~\date{}}
	\fancyhead[C]{}
	\fancyhead[R]{\textbf{Solutions}}
	\setlength{\headsep}{14pt}
}
%%%%%%%%%%%%%%%%%%%%%%%%%%%%%%%%%%%%%%%%%%

%%Next page's header%%%%%%%%%%%%%%%%%%%%%%%
\setlist{nolistsep}
\lhead{\classNum{}(\classSec{})} 
\chead{\textbf{Solutions}} 
\cfoot{\thepage{}~of~\pageref{LastPage}} 
\setlength{\headsep}{5pt} 
%%%%%%%%%%%%%%%%%%%%%%%%%%%%%%%%%%%%%%%%%%%

%%Various settings; see the preamble for other settings%%%%%
%\setenumerate[0]{leftmargin=12pt,itemindent=0pt,labelwidth=0pt}
\setlist{topsep=.5pt}
\renewcommand{\labelenumi}{\textbf{\arabic{enumi}.}}
\renewcommand{\labelenumii}{\textbf{\alph{enumii}.}}
\thispagestyle{firststyle}
%%%%%%%%%%%%%%%%%%%%%%%%%%%%%%%%%%%%%%%%%%

\newcommand{\xmax}{0}
\newcommand{\xmin}{-\xmax}
\newcommand{\xstep}{0}
\newcommand{\ymax}{0}
\newcommand{\ymin}{-\ymax}
\newcommand{\ystep}{0} 
\newcommand{\dspace}{\vspace{1cm}}
\newcommand{\xa}{0}
\newcommand{\xb}{0}
\newcommand{\ya}{1}
\newcommand{\yb}{1}
\begin{enumerate}[resume]
\item \vspace{0cm}Note: this are the same matrices from 5.1 exercise 1.

\begin{enumerate}[resume]
\item Oops! I did this one in class. $\det\left(A-tI_{2}\right)=\det\begin{bmatrix}\begin{array}{rr}
1-t & 5\\
2 & 4-t
\end{array}\end{bmatrix}=(1-t)(4-t)-10=t^{2}-5t-6=(t-6)(t+1)=0$ when $t=6,-1$ so $A$ has eigenvalues $\lambda=-1,6$. 

\begin{enumerate}[resume]
\item []$E(-1):$ $N(A+I_{2})=N\left(\begin{bmatrix}\begin{array}{rr}
1+1 & 5\\
2 & 4+1
\end{array}\end{bmatrix}\right)=N\left(\begin{bmatrix}\begin{array}{rr}
2 & 5\\
2 & 5
\end{array}\end{bmatrix}\right)$ and $\begin{bmatrix}\begin{array}{rr}
2 & 5\\
2 & 5
\end{array}\end{bmatrix}\rightsquigarrow\begin{bmatrix}\begin{array}{rr}
2 & 5\\
0 & 0
\end{array}\end{bmatrix}$ so $2x_{1}=-5x_{2}\Rightarrow\mathbf{v}_{1}=\begin{bmatrix}\begin{array}{r}
5\\
-2
\end{array}\end{bmatrix}$ gives a basis for $E(-1)$.
\item []$E(6):$ $N(A-6I_{2})=N\left(\begin{bmatrix}\begin{array}{rr}
1-6 & 5\\
2 & 4-6
\end{array}\end{bmatrix}\right)=N\left(\begin{bmatrix}\begin{array}{rr}
-5 & 5\\
2 & -2
\end{array}\end{bmatrix}\right)$ and $\begin{bmatrix}\begin{array}{rr}
-5 & 5\\
2 & -2
\end{array}\end{bmatrix}\rightsquigarrow\begin{bmatrix}\begin{array}{rr}
1 & -1\\
0 & 0
\end{array}\end{bmatrix}$ so $x_{1}=x_{2}\Rightarrow\mathbf{v}_{2}=\begin{bmatrix}\begin{array}{r}
1\\
1
\end{array}\end{bmatrix}$ gives a basis for $E(6)$.
\end{enumerate}
\item $\det\left(A-tI_{2}\right)=\det\begin{bmatrix}\begin{array}{rr}
0-t & 1\\
1 & 0-t
\end{array}\end{bmatrix}=(-t)(-t)-1=t^{2}-1=(t-1)(t+1)=0$ when $t=1,-1$ so $A$ has eigenvalues $\lambda=-1,1$. 

\begin{enumerate}[resume]
\item []$E(-1):$ $N(A+I_{2})=N\left(\begin{bmatrix}\begin{array}{rr}
1 & 1\\
1 & 1
\end{array}\end{bmatrix}\right)$ and $\begin{bmatrix}\begin{array}{rr}
1 & 1\\
1 & 1
\end{array}\end{bmatrix}\rightsquigarrow\begin{bmatrix}\begin{array}{rr}
1 & 1\\
0 & 0
\end{array}\end{bmatrix}$ so $x_{1}=-x_{2}=0\Rightarrow\mathbf{v}_{1}=\begin{bmatrix}\begin{array}{r}
1\\
-1
\end{array}\end{bmatrix}$ gives a basis for $E(-1)$.
\item []$E(1):$ $N(A-I_{2})=N\left(\begin{bmatrix}\begin{array}{rr}
-1 & 1\\
1 & -1
\end{array}\end{bmatrix}\right)$ and $\begin{bmatrix}\begin{array}{rr}
-1 & 1\\
1 & -1
\end{array}\end{bmatrix}\rightsquigarrow\begin{bmatrix}\begin{array}{rr}
1 & -1\\
0 & 0
\end{array}\end{bmatrix}$ so $x_{1}=x_{2}\Rightarrow\mathbf{v}_{2}=\begin{bmatrix}\begin{array}{r}
1\\
1
\end{array}\end{bmatrix}$ gives a basis for $E(1)$.
\end{enumerate}
\item $\det\left(A-tI_{2}\right)=\det\begin{bmatrix}\begin{array}{rr}
10-t & -6\\
18 & -11-t
\end{array}\end{bmatrix}=(10-t)(-11-t)+108=t^{2}+t-2=(t-1)(t+2)=0$ when $t=1,-2$ so $A$ has eigenvalues $\lambda=-2,1$. 

\begin{enumerate}[resume]
\item []$E(-2):$ $N(A+2I_{2})=N\left(\begin{bmatrix}\begin{array}{rr}
12 & -6\\
18 & -9
\end{array}\end{bmatrix}\right)$ and $\begin{bmatrix}\begin{array}{rr}
12 & -6\\
18 & -9
\end{array}\end{bmatrix}\rightsquigarrow\begin{bmatrix}\begin{array}{rr}
1 & -\nicefrac{1}{2}\\
0 & 0
\end{array}\end{bmatrix}$ so $x_{1}=\nicefrac{1}{2}x_{2}=0\Rightarrow\mathbf{v}_{1}=\begin{bmatrix}\begin{array}{r}
1\\
\nicefrac{1}{2}
\end{array}\end{bmatrix}$ gives a basis for $E(-2)$.
\item []$E(1):$ $N(A-I_{2})=N\left(\begin{bmatrix}\begin{array}{rr}
-1 & 1\\
1 & -1
\end{array}\end{bmatrix}\right)$ and $\begin{bmatrix}\begin{array}{rr}
-1 & 1\\
1 & -1
\end{array}\end{bmatrix}\rightsquigarrow\begin{bmatrix}\begin{array}{rr}
1 & -1\\
0 & 0
\end{array}\end{bmatrix}$ so $x_{1}=x_{2}\Rightarrow\mathbf{v}_{2}=\begin{bmatrix}\begin{array}{r}
1\\
1
\end{array}\end{bmatrix}$ gives a basis for $E(1)$.
\end{enumerate}
\item [\textbf{j.}]$\det\left(A-tI_{3}\right)=\det\begin{bmatrix}\begin{array}{rrr}
2-t & 0 & 1\\
0 & 1-t & 2\\
0 & 0 & 1-t
\end{array}\end{bmatrix}=(1-t)\det\left(\begin{bmatrix}\begin{array}{rr}
2-t & 1\\
0 & 1-t
\end{array}\end{bmatrix}\right)=(1-t)(2-t)(1-t)=0$ when $t=1,2$ so $A$ has eigenvalues $\lambda=1,2$. 

\begin{enumerate}[resume]
\item []$E(1):$ $N(A-I_{3})=N\left(\begin{bmatrix}\begin{array}{rrr}
1 & 0 & 1\\
0 & 0 & 2\\
0 & 0 & 0
\end{array}\end{bmatrix}\right)$ and $\begin{bmatrix}\begin{array}{rrr}
1 & 0 & 1\\
0 & 0 & 2\\
0 & 0 & 0
\end{array}\end{bmatrix}\rightsquigarrow\begin{bmatrix}\begin{array}{rrr}
1 & 0 & 0\\
0 & 0 & 1\\
0 & 0 & 0
\end{array}\end{bmatrix}$ so $x_{1}=x_{3}=0\Rightarrow\mathbf{v}_{1}=\begin{bmatrix}\begin{array}{r}
0\\
1\\
0
\end{array}\end{bmatrix}$
\item []$E(2):$ $N(A-2I_{3})=N\left(\begin{bmatrix}\begin{array}{rrr}
0 & 0 & 1\\
0 & -1 & 2\\
0 & 0 & -1
\end{array}\end{bmatrix}\right)$ and $\begin{bmatrix}\begin{array}{rrr}
0 & 0 & 1\\
0 & -1 & 2\\
0 & 0 & -1
\end{array}\end{bmatrix}\rightsquigarrow\begin{bmatrix}\begin{array}{rrr}
0 & 1 & 0\\
0 & 0 & 1\\
0 & 0 & 0
\end{array}\end{bmatrix}$ so $x_{2}=x_{3}=0\Rightarrow\mathbf{v}_{2}=\begin{bmatrix}\begin{array}{r}
1\\
0\\
0
\end{array}\end{bmatrix}$
\end{enumerate}
\item [\textbf{m.}]$\det\left(A-tI_{3}\right)=\det\begin{bmatrix}\begin{array}{rrr}
1-t & -1 & 4\\
3 & 2-t & -1\\
2 & 1 & -1-t
\end{array}\end{bmatrix}=(t-1)(t+2)(t-3)=0$ when $t=1,-2,3$ so $A$ has eigenvalues $\lambda=1,-2,3$ (done
in class). 

\begin{enumerate}[resume]
\item []$E(-2):$ $N(A+2I_{3})=N\left(\begin{bmatrix}\begin{array}{rrr}
3 & -1 & 4\\
3 & 4 & -1\\
2 & 1 & 1
\end{array}\end{bmatrix}\right)$ and $\begin{bmatrix}\begin{array}{rrr}
3 & -1 & 4\\
3 & 4 & -1\\
2 & 1 & 1
\end{array}\end{bmatrix}\rightsquigarrow\begin{bmatrix}\begin{array}{rrr}
1 & 0 & 1\\
0 & 1 & -1\\
0 & 0 & 0
\end{array}\end{bmatrix}$ so $x_{1}=-x_{3}\mbox{ and }x_{2}=x_{3}\Rightarrow\mathbf{v}_{1}=\begin{bmatrix}\begin{array}{r}
-1\\
1\\
1
\end{array}\end{bmatrix}$
\item []$E(1):$ $N(A-I_{3})=N\left(\begin{bmatrix}\begin{array}{rrr}
0 & -1 & 4\\
3 & 1 & -1\\
2 & 1 & -2
\end{array}\end{bmatrix}\right)$ and $\begin{bmatrix}\begin{array}{rrr}
0 & -1 & 4\\
3 & 1 & -1\\
2 & 1 & -2
\end{array}\end{bmatrix}\rightsquigarrow\begin{bmatrix}\begin{array}{rrr}
1 & 0 & 1\\
0 & 1 & -4\\
0 & 0 & 0
\end{array}\end{bmatrix}$ so $x_{1}=-x_{3}\mbox{ and }x_{2}=4x_{3}\Rightarrow\mathbf{v}_{2}=\begin{bmatrix}\begin{array}{r}
-1\\
4\\
1
\end{array}\end{bmatrix}$
\item []$E(3):$ $N(A-3I_{3})=N\left(\begin{bmatrix}\begin{array}{rrr}
-2 & -1 & 4\\
3 & -1 & -1\\
2 & 1 & -4
\end{array}\end{bmatrix}\right)$ and $\begin{bmatrix}\begin{array}{rrr}
-2 & -1 & 4\\
3 & -1 & -1\\
2 & 1 & -4
\end{array}\end{bmatrix}\rightsquigarrow\begin{bmatrix}\begin{array}{rrr}
1 & 0 & -1\\
0 & 1 & -2\\
0 & 0 & 0
\end{array}\end{bmatrix}$ so $x_{2}=x_{3}\mbox{ and }x_{2}=2x_{3}\Rightarrow\mathbf{v}_{3}=\begin{bmatrix}\begin{array}{r}
1\\
2\\
1
\end{array}\end{bmatrix}$
\end{enumerate}
\end{enumerate}
\end{enumerate}

\end{document}
